% !TeX program = pdflatex
% !BIB program = biber
%
% Scientific protocol of the practical work, ready to paste into Overleaf.
% Compiles with pdflatex + biber (Overleaf default since 2017).
%
% To use in Overleaf:
%   1. Create a new project; upload Protocol.tex and References.bib.
%   2. Project settings -> compiler: pdfLaTeX (default).
%   3. Recompile; biber runs automatically for biblatex.

\documentclass[11pt,a4paper]{article}

\usepackage[utf8]{inputenc}
\usepackage[T1]{fontenc}
\usepackage{lmodern}
\usepackage{microtype}
\usepackage{amsmath,amssymb}
\usepackage{graphicx}
\usepackage{booktabs}
\usepackage{xcolor}
\usepackage[a4paper, margin=2.5cm]{geometry}
\usepackage{hyperref}
\hypersetup{colorlinks=true,
            linkcolor=blue!45!black,
            citecolor=blue!45!black,
            urlcolor=blue!45!black}
\usepackage[backend=biber,
            style=authoryear-comp,
            sorting=nyt,
            maxcitenames=2,
            maxbibnames=10,
            uniquelist=false]{biblatex}
\addbibresource{References.bib}

\title{Practical Work Protocol:\\[0.3em]
CT-to-PET Translation via Brownian Bridge Diffusion on autoPET Dataset}
\author{Kacper Zaleski\\Johannes Kepler University Linz}

\date{\today}

\begin{document}
\maketitle

\begin{abstract}
This protocol documents the end-to-end reimplementation of CPDM, a 
domain-knowledge-guided diffusion model for CT-to-PET translation 
\parencite{nguyen2025cpdm}, on the \texttt{autoPET} whole-body
multi-tracer PET/CT dataset \parencite{gatidis2022autopet, jeblick2024psma}. The pipeline comprises (1)
an attention-map UNet that learns regions of high PET uptake from CT alone, (2) 
a VQGAN encoder/decoder providing a $4 \times 16 \times 16$ latent space, and (3) 
a Brownian Bridge diffusion process \parencite{li2023bbdm} operating in that latent 
space, conditioned on the attention map and a CT-derived 511\,keV attenuation map. 
All four training stages were executed CPU-only on an Intel i7-8650U 
(8 cores, 15\,GB RAM). The model reaches a best validation LPIPS of 0.166 and 
PSNR of 41.7\,dB on \texttt{autoPET} brain-region slices, beating three trivial 
baselines on PSNR by 2--16\,dB and confirming that the model produces structurally
meaningful PET despite a background-dominated data distribution.
\end{abstract}

\tableofcontents
\bigskip

\section{Introduction}

Computed tomography (CT), introduced by \textcite{hounsfield1973ct}, reconstructs 
cross-sectional anatomical images from rotating X-ray projections and remains the 
workhorse modality for structural imaging across nearly every medical specialty. 
Positron emission tomography (PET), developed shortly after 
by \textcite{phelps1975pet}, instead images the spatial distribution of an 
injected radiotracer (most commonly $^{18}$F-fluorodeoxyglucose, FDG) via 
coincident-photon detection, yielding \emph{functional} information about 
tissue metabolism that CT alone cannot resolve \parencite{gambhir2002pet}. 
The two modalities are typically acquired together on hybrid PET/CT scanners 
introduced by \textcite{beyer2000petct}, where the CT additionally provides 
the attenuation correction needed to quantify the PET signal. PET, however, 
carries costs that CT does not: an additional ionising-radiation dose from 
the radiotracer on top of the CT component, longer acquisition times, the need 
for an on-site cyclotron or tracer supply chain, and substantially higher per-study 
costs. Repeated CT studies themselves now constitute a substantial source of 
population-level radiation exposure \parencite{brennerhall2007ctdose}, making 
any reduction in additional PET dose clinically relevant. The ability to 
synthesise PET from CT even imperfectly, therefore has clinical value as a 
triage tool, as a data-augmentation aid for downstream tasks like tumour 
segmentation, and as a research vehicle for understanding what physical and 
anatomical information CT can convey about radiotracer uptake.

This project reimplements \textbf{CPDM} (CT to PET Diffusion Model) \parencite{nguyen2025cpdm}, 
a domain-knowledge-guided diffusion approach proposed at WACV 2025 using the authors' 
large-scale curated CT-PET dataset. The approach is reproduced on the publicly 
available \texttt{autoPET} dataset \parencite{gatidis2022autopet, jeblick2024psma}, a whole-body
multi-tracer (FDG and PSMA) PET/CT dataset with tumour segmentation labels, and every adaptation required
to fit the original paper's design into the available compute budget and the dataset's 
properties (HU-normalised CT, varied scanner conditions, partially imbalanced 
tumour distributions) is documented.

The aim of the practical work is a faithful reproduction of the CPDM architecture 
and training pipeline on \texttt{autoPET}, with all paper-required 
components: VQGAN, Brownian Bridge denoiser, and the two domain-specific 
conditioning signals.

\section{Theoretical Background}

\subsection{Denoising diffusion probabilistic models}

Diffusion models \parencite{sohldickstein2015deep, ho2020denoising, song2021scorebased} learn a generative process by inverting a fixed Gaussian noising process. Given a data distribution $q(x_0)$, the forward process produces a sequence of progressively noisier latents $x_1, \ldots, x_T$ via the Markov chain
\begin{equation}
q(x_t \mid x_{t-1}) = \mathcal{N}\!\left(x_t;\, \sqrt{1 - \beta_t}\, x_{t-1},\, \beta_t\, I\right)
\end{equation}
with a variance schedule $\beta_1 < \cdots < \beta_T$. A neural network 
$\varepsilon_\theta(x_t, t)$ is trained to predict the noise added at each 
step, equivalently learning the score $\nabla_{x_t} \log p(x_t)$ \parencite{song2021scorebased}. Sampling 
proceeds by iteratively denoising from $x_T \sim \mathcal{N}(0, I)$.

The original DDPM \parencite{ho2020denoising} uses a UNet denoiser 
\parencite{ronneberger2015unet} adapted with time-step embeddings, 
group normalisation, and attention layers at low spatial resolutions. 
Subsequent work demonstrated that diffusion models in pixel space at 
high resolutions are computationally expensive. Latent diffusion models (LDMs) 
\parencite{rombach2022ldm} address this by training the diffusion process in the 
latent space of a pretrained autoencoder, dramatically reducing compute while 
preserving sample quality. The autoencoder is typically either a VAE 
\parencite{kingma2014vae} with a continuous, KL-regularised latent or 
a VQGAN \parencite{esser2021vqgan}; CPDM, and therefore this 
work, uses the latter. VQGAN itself builds on VQ-VAE \parencite{oord2017vqvae}, 
which replaces the VAE's continuous Gaussian latent with a discrete codebook of 
learned code vectors, and adds a perceptual loss together with an adversarial 
discriminator to obtain sharper reconstructions.

\subsection{Brownian Bridge diffusion for image-to-image translation}
\label{sec:bbdm}

Standard DDPM samples are unconditional. For paired image-to-image translation (e.g.\ CT $\to$ PET), 
one wants a diffusion process whose endpoint distributions are \emph{both} data 
distributions, not one data and one Gaussian noise. BBDM \parencite{li2023bbdm} 
proposes a diffusion bridge between source and target image distributions using a 
Brownian Bridge stochastic process: a Brownian motion conditioned on its endpoints. 
The forward process at time $t$ interpolates between $x_0$ (target) and $y$ (source) 
according to
\begin{equation}
x_t \;=\; (1 - m_t)\, x_0 \;+\; m_t\, y \;+\; \sigma_t\, \varepsilon_t,
\label{eq:bbdm-fwd}
\end{equation}
where $m_t$ is a monotone schedule from $0$ to $1$ and $\sigma_t$ controls 
injected noise. At $t = 0$ the latent equals the target; at $t = T$ it equals 
the source. The reverse process is trained to predict the gradient of the bridge, 
and sampling iteratively transforms the source into the target. Compared to 
the older approach of diffusion conditioning, where the source image
is simply concatenated to the noisy target as a context channel, BBDM has a 
more principled 
forward process for image-to-image tasks and empirically produces sharper 
results on benchmarks like CelebA-HQ-edge and Cityscapes \parencite{li2023bbdm}.

\subsection{CPDM: domain-knowledge-guided BB diffusion for CT-to-PET}
\label{sec:cpdm-theory}

The CPDM paper \parencite{nguyen2025cpdm} adapts BBDM to medical CT-to-PET 
translation with two domain-specific conditioning signals concatenated as 
additional input channels to the BB denoiser:

\begin{enumerate}
  \item \textbf{Attention map} ($M_T$): a binary spatial mask indicating regions 
  of expected high PET uptake. Produced by a separately trained UNet that learns 
  the mapping $\text{CT} \to (\text{PET} > p_{75} \;+\; \text{binary closing})$. 
  Provides the denoiser with a hard spatial prior about \emph{where} to place 
  bright regions.
  \item \textbf{Attenuation map} ($M_\mu$): a closed-form per-pixel transform of 
  CT Hounsfield units to the linear attenuation coefficient at 511\,keV the energy 
  of the annihilation photons whose detection PET measures. Approximated by a 
  piecewise-linear function calibrated to CT scanner kVp. Encodes the physics 
  linking CT density and PET signal correction.
\end{enumerate}

The paper's ablation (their Table~4) reports that removing $M_T$ increases LPIPS 
by $\sim 12\%$ and removing $M_\mu$ by $\sim 14\%$, supporting the claim that 
both signals contribute meaningfully.

CPDM operates in the latent space of a pretrained VQGAN trained on the same 
paired data, following the LDM paradigm. The denoising UNet has input 
channels $3 \cdot z_{\text{channels}}$ (target latent + two condition 
contexts, each remapped via a \texttt{SpatialRescaler} to $z_{\text{channels}}$), 
matching the $9 = 3 \cdot 3$ design used in the released code template.

\subsection{Evaluation metrics}
\label{sec:eval-metrics}

A standard image-quality metric suite for translation tasks includes:
\begin{itemize}
  \item \textbf{Mean Absolute Error (MAE)}: per-pixel $\lVert x - \hat{x} \rVert_1$. 
  Sensitive to intensity bias and noise; widely reported but trivially minimised by 
  predicting the per-pixel mean of the training distribution.
  \item \textbf{Peak Signal-to-Noise Ratio (PSNR)}: $10 \cdot \log_{10}(R^2 / \text{MSE})$ 
  in dB, where $R$ is the dynamic range. Rewards low overall error; less easily 
  fooled by mean-image predictors than MAE alone.
  \item \textbf{Structural Similarity Index (SSIM)}: \parencite{wang2004ssim} 
  captures luminance, contrast, and structural correlation in local windows. 
  Penalises loss of structure; more clinically meaningful than per-pixel 
  metrics for medical images.
  \item \textbf{Learned Perceptual Image Patch Similarity (LPIPS)}: \parencite{zhang2018lpips} 
  distance in the activation space of a pretrained CNN (here, VGG). Correlates 
  better than the above with human judgments of perceptual similarity for natural 
  images; less well validated for medical images but widely adopted.
\end{itemize}

The CPDM paper \parencite{nguyen2025cpdm} reports all four; the same suite is 
replicated here to enable direct (paper-comparable) numbers.

\section{Dataset}

\subsection{\texttt{autoPET}}

This work uses the \texttt{autoPET III} release, which combines two publicly
distributed, tumour-lesion-annotated whole-body PET/CT cohorts acquired with two
different radiotracers: the FDG-PET/CT cohort of \textcite{gatidis2022autopet}
(1{,}014 studies, 900 patients; UKT T\"ubingen) and the PSMA-PET/CT cohort of
\textcite{jeblick2024psma} (LMU Munich). The combined set used here totals
\textbf{1{,}614} studies (1{,}014 FDG + 600 PSMA), distributed via the
\texttt{autoPET} challenge held at MICCAI 2022--2024. Each study contains a paired
CT and PET volume in NIfTI format together with manually annotated tumour
segmentation masks. Because FDG (metabolic) and PSMA (prostate-specific membrane
antigen) label different biological targets, the combined set spans a wider range
of uptake patterns than an FDG-only set would.

For this practical work the official 80/10/10 split (\texttt{splits\_80\_10\_10.json}
shipped with the dataset) is used: 1{,}291 train / 161 val / 162 test studies, whose
tracer composition is shown in Table~\ref{tab:splits}. Only the \textbf{first 20\,\%
of slices} per volume corresponding approximately to the head/brain region are
retained, to keep the dataset size manageable on CPU.

\begin{table}[h]
\centering
\begin{tabular}{lrrrr}
\toprule
Split & FDG & PSMA & Studies & $64 \times 64$ brain slices \\
\midrule
train & 818 & 473 & 1{,}291 & 86{,}197 \\
val   & 102 & 59  & 161    & 10{,}464 \\
test  & 94  & 68  & 162    & 11{,}330 \\
\midrule
total & 1{,}014 & 600 & 1{,}614 & 107{,}991 \\
\bottomrule
\end{tabular}
\caption{\texttt{autoPET III} study counts and brain-region slice counts after the
80/10/10 split, broken down by radiotracer.}
\label{tab:splits}
\end{table}

This brain-only restriction is a known deviation from the paper, which uses 
scans from top of head to the upper thigh region. It restricts the diversity of 
uptake patterns (no liver, bladder, or heart background uptake) and inflates the proportion 
of near-background pixels.

\subsection{Preprocessing}

Preprocessing (\texttt{preprocess\_autopet.py}) performs the following steps per volume:
\begin{enumerate}
  \item \textbf{Extract} the first 20\% of axial slices.
  \item \textbf{Resize} each 2D slice to $64 \times 64$. 
  \item \textbf{Normalise} to $[-1, 1]$. CT is clipped to HU $\in [-1000, 3071]$ and 
  PET to SUV $\in [0, 32]$; each clipped value $v$ in $[a, b]$ is then mapped affinely 
  as $\tilde v = 2(v-a)/(b-a) - 1$. So CT HU $-1000 \mapsto -1$, $3071 \mapsto +1$; 
  PET SUV $0 \mapsto -1$, $32 \mapsto +1$. This puts both modalities on a common, 
  symmetric range for the VQGAN and diffusion model.
  \item \textbf{Save} each slice as a \texttt{float32} NumPy file
\end{enumerate}

\section{Method}

\subsection{Stage 1: Attention-map UNet}

\paragraph{Architecture.} A 2D UNet \parencite{ronneberger2015unet} implemented via \texttt{segmentation\_models\_pytorch} \parencite{iakubovskii2019smp} with a ResNet34 \parencite{he2016resnet} encoder (ImageNet-pretrained weights averaged into the single-channel input) and a concurrent spatial-and-channel squeeze-and-excitation (scSE) decoder \parencite{roy2018scse}.

\paragraph{Target.} The CPDM paper specifies the attention map as a binary mask of high-uptake regions derived from the paired PET, but leaves the exact threshold and post-processing implementation-defined. We follow a simple, reproducible recipe consistent with that description:
\begin{enumerate}
  \item Compute the per-slice 75th percentile $\tau$ of the (normalised) PET intensities; mark every pixel with intensity $> \tau$ as foreground. This isolates the brightest quartile of pixels, which on PET corresponds to elevated tracer uptake (FDG or PSMA).
  \item Apply binary morphological closing (dilation followed by erosion) to the resulting mask using a $3 \times 3$ all-ones structuring element. The structuring element is the small neighbourhood the morphological operation slides over each pixel; a $3 \times 3$ square is the smallest isotropic 8-connected neighbourhood and is the standard default. Closing fills single-pixel holes and joins near-adjacent foreground blobs into contiguous regions, producing a cleaner training target than raw thresholding.
\end{enumerate}
The mask is computed on the fly from PET inside \texttt{AttentionMapDataset}; CT is the network input and PET is held out during attention-UNet training, so at CPDM training time the attention map is derived from CT alone.

\paragraph{Loss.} $0.7 \cdot \text{DiceLoss} + 0.3 \cdot \text{BCEWithLogitsLoss}$. 
The Dice term \parencite{milletari2016vnet} dominates to focus optimisation on 
overlap with the foreground; the BCE term stabilises gradients in the early 
training when predictions are near-uniform.

\paragraph{Training.} Adam optimiser \parencite{kingma2015adam}, learning rate 
$10^{-3}$, \texttt{CosineAnnealingLR} schedule, batch size 8, 3 epochs $\times$ 
200 batches-per-epoch limit. 

\subsection{Stage 2: Attention-map export}

After training, \texttt{export\_attention\_maps.py} runs the trained UNet 
over every CT slice in all three splits at batch size 32, writing the sigmoid 
probability per slice as a \texttt{float32} array of shape $(64, 64)$ to 
\texttt{data/processed/\{split\}/AttentionMaps/<slice>.npy}. This is the 
exact format the CPDM denoiser reads at training time, with a hard threshold 
of 0.5 applied at load. Total writes: 108k maps.

\subsection{Stage 3: VQGAN encoder/decoder}

\paragraph{Architecture.} The VQGAN \parencite{esser2021vqgan} from the original CPDM codebase, configured for single-channel medical imaging: base channels = 128, channel multipliers $(1, 2, 4)$, two residual blocks per level, no attention layers, latent channels $z = 4$, codebook size 8192. Spatial compression $64 \to 16$ ($\times 4$ downsampling). Total trainable parameters: 55.3\,M. The encoder--quantiser--decoder structure follows \textcite{oord2017vqvae} for VQ-VAE.

\paragraph{Training.} $L_1$ reconstruction loss plus the standard commitment/codebook 
loss; the \textbf{perceptual loss} \parencite{zhang2018lpips} and 
\textbf{GAN discriminator} \parencite{goodfellow2014gan, esser2021vqgan} used 
in the original VQGAN are not used here.

\paragraph{Schedule.} Adam optimiser, learning rate $10^{-4}$, betas $(0.5, 0.9)$. Each batch concatenates CT and PET along the batch dimension so the encoder learns both modalities equally. CPU-friendly settings: batch size 4 (effective 8 after CT$\|$PET concat), 5 epochs $\times$ 500 batches-per-epoch limit. Final val loss = 0.0185 (rec = 0.0168, codebook = 0.0017) at epoch 1; further training did not improve metrics meaningfully.

\subsection{Stage 4: CPDM (Brownian Bridge in latent space)}

\paragraph{Architecture.} \texttt{CT2PETDiffusionModel} includes:
\begin{itemize}
  \item Frozen VQGAN (loaded from the Stage 3 checkpoint) encoding 
  $1 \times 64 \times 64 \to 4 \times 16 \times 16$.
  \item Two \texttt{SpatialRescaler} condition stages, one each for 
  attention and attenuation maps. Each performs $n_{\text{stages}} = 2$ 
  bilinear downsamplings with multiplier 0.5 ($64 \to 32 \to 16$) followed by a 
  learned $1 \times 1$ \texttt{Conv2d} remapping $1 \to 4$ channels 
  (matching $z_{\text{channels}}$).
  \item A \textbf{UNet denoiser} \parencite{ronneberger2015unet, ho2020denoising} 
  adapted from the OpenAI guided-diffusion codebase. Configuration: 
  \texttt{image\_size} = 16 (operates on the latent), \texttt{model\_channels} = 128,
  $\texttt{channel\_mult} = (1, 2, 3, 4)$, two residual blocks per level, attention 
  at resolutions $\{8, 4\}$ with \texttt{num\_heads} = 8, scale-shift normalisation. 
  Input channels = $4\,(x_t) + 4\,(\text{attention context}) + 4\,(\text{attenuation context}) = 12$.
  Output channels = 4.
\end{itemize}

\paragraph{Conditioning signals.}
\begin{itemize}
  \item \emph{Attention map}: pre-exported in Stage 2; loaded by filename from 
  \texttt{data/processed/\{split\}/AttentionMaps/} and thresholded at 0.5 inside 
  \texttt{CT2PETDiffusionModel.get\_attention\_map}.
  \item \emph{Attenuation map}: closed-form transform of CT HU to 511\,keV linear
  attenuation coefficient via the piecewise-linear approximation from PET
  attenuation correction literature, parameterised by scanner kVp (set to
  kVp = 140 here). The transform is applied on the fly inside
  \texttt{CT2PETDiffusionModel.get\_attenuation\_map} and expects HU-valued
  input; the preceding line in that function reverses the $[-1, 1]$
  normalisation applied by \texttt{normalize\_ct} so the piecewise-linear
  lookup receives HU as intended.
\end{itemize}

\paragraph{Diffusion process} Brownian Bridge \parencite{li2023bbdm} with 
1000 forward timesteps and 200 reverse sampling timesteps, monotone schedule 
(\texttt{mt\_type = 'linear'}), maximum variance 1.0, $\eta = 1.0$. 
Training objective: \texttt{grad} (the bridge gradient), loss type 
$L_1$. Skip sampling is enabled for the accelerated DDIM-style reverse 
process \parencite{song2022ddim}.

\paragraph{Training schedule} Adam optimiser, learning rate $10^{-4}$, 
no weight decay. \texttt{ReduceLROnPlateau} scheduler 
\parencite{robbins1951stochastic} with patience 3000 
iterations and factor 0.5. The \texttt{CPDMRunner} overrides 
\texttt{validation\_epoch} to (i) compute the standard val loss, 
(ii) generate paper metrics on test batches via the full 
200-step BB reverse process, and (iii) check an early-stopping 
criterion (patience = 15 epochs on \texttt{val\_epoch/loss}, 
$\Delta_{\min} = 5 \times 10^{-4}$). The patience is intentionally 
larger than the LR scheduler's effective patience so at least one 
LR drop is observed before stoping.

\section{Implementation Details}

\subsection{Compute environment}

Intel i7-8650U, 8 cores @ 1.9\,GHz, 15\,GB RAM. 

\subsection{Software stack}

PyTorch 2.9.0 \parencite{paszke2019pytorch}, PyTorch Lightning 2.6 \parencite{falcon2019lightning}, \texttt{lpips} 0.1.4, \texttt{scikit-image} 0.26 \parencite{walt2014scikitimage}, \texttt{wandb} 0.25 \parencite{biewald2020wandb}, \texttt{nibabel} 5.3, \texttt{scipy} 1.17. Full environment specification in \texttt{pyproject.toml}.

\subsection{Documented adaptations from the paper}

\begin{table}[h]
\centering
\small
\begin{tabular}{p{0.5cm}p{6.2cm}p{6.5cm}}
\toprule
\# & Adaptation & Reason \\
\midrule
1 & Input size $64 \times 64$ vs.\ paper's $256 \times 256$ & CPU budget. \\
2 & Brain region only (first 20\% slices) vs\ almost full body & CPU budget. \\
3 & Attention UNet: ResNet34 + $0.7\cdot\text{Dice} + 0.3\cdot\text{BCE}$ vs.\ paper's ResNet50 + pure Dice & Smaller backbone for CPU; BCE blend for early-training stability. \\
4 & Wrote \texttt{train\_vqgan.py} & Paper expects external VQGAN training via \texttt{taming-transformers} \parencite{esser2021vqgan}. \\
\bottomrule
\end{tabular}
\caption{Adaptations to the original CPDM pipeline.}
\label{tab:adaptations}
\end{table}

\section{Experiments and Results}

\subsection{Stage 1 --- Attention UNet}

After 3 epochs $\times$ 200 batches at batch size 8 on CPU (${\sim}6$ minutes total):
\begin{itemize}
  \item \texttt{val\_dataset\_iou} $= \mathbf{0.7050}$,
  \item \texttt{val\_per\_image\_iou} $= 0.7079$,
  \item \texttt{val\_f1\_score} $= \mathbf{0.8266}$,
  \item train epoch loss $= 0.246$.
\end{itemize}
The model learns the mapping $\text{CT} \to$ high-uptake mask reliably. Qualitative samples in \texttt{report\_out/02\_attention\_unet\_predictions.png} show the predicted masks aligning with the PET-derived targets on representative val slices.

\subsection{Stage 2 --- Attention-map export}

108\,k attention maps written ($86{,}197 + 10{,}464 + 11{,}330$). On-disk shape $(64, 64)$ \texttt{float32} with values in $[0, 1]$; sample mean $\approx 0.32$, range $= [0, 1]$. The export format is byte-compatible with \texttt{CT2PETDiffusionModel.get\_attention\_map}'s loader.

\subsection{Stage 3 --- VQGAN}

After Stage 3 training, validation loss reaches $0.0185$ at epoch 1 with $\texttt{val\_rec\_loss} = 0.0168$ and $\texttt{val\_codebook\_loss} = 0.0017$. Reconstructions on val examples (figure \texttt{04\_vqgan\_reconstructions.png}) preserve the dominant intensity structure but smooth fine detail --- expected given the $4\times$ spatial compression and absence of perceptual / GAN losses.

\subsection{Stage 4 --- CPDM training trajectory}

Two long-run launches and one resume covering 10 completed epochs (250 iterations each, batch size 8):

\begin{table}[h]
\centering
\small
\begin{tabular}{rrrrrr}
\toprule
Epoch & \texttt{val\_epoch/loss} & LPIPS$\downarrow$ & MAE$\downarrow$ & SSIM$\uparrow$ & PSNR$\uparrow$ (dB) \\
\midrule
0 & 0.02184          & 0.186 & 0.0055      & 0.880          & 41.35 \\
1 & 0.02003          & 0.176 & \textbf{0.0053} & \textbf{0.890} & \textbf{41.70} \\
2 & 0.02053          & 0.174 & 0.0055      & 0.876          & 41.41 \\
3 & 0.02046          & 0.171 & 0.0054      & 0.879          & 41.55 \\
4 & 0.02103          & 0.200 & 0.0059      & 0.880          & 40.94 \\
5 & 0.02080          & 0.192 & 0.0057      & 0.874          & 41.16 \\
6 & 0.01977          & \textbf{0.166} & 0.0056 & 0.875       & 41.44 \\
7 & \textbf{0.01394} & 0.176 & 0.0055      & 0.879          & 41.45 \\
8 & 0.01712          & 0.200 & 0.0059      & 0.868          & 40.77 \\
9 & ---              & 0.174 & 0.0055      & 0.878          & 41.46 \\
\bottomrule
\end{tabular}
\caption{Per-epoch metrics on \texttt{autoPET} validation slices. Bold entries are best on the column.}
\label{tab:cpdm-trajectory}
\end{table}

Best \texttt{val\_epoch/loss} of $0.01394$ at epoch 7. Metrics oscillate within a narrow band across all 10 epochs (LPIPS $\pm 0.03$, MAE $\pm 0.0006$, SSIM $\pm 0.02$, PSNR $\pm 0.9$\,dB). The \texttt{ReduceLROnPlateau} scheduler's effective patience ($\approx 12$ epochs at 250 iter/epoch) had not triggered at the manual stop. Training was halted at epoch 11 because all four metrics had reached a stable plateau and the visual quality was no longer improving qualitatively.

\begin{figure}[h]
  \centering
  \includegraphics[width=\linewidth]{../report_out/05_cpdm_training_curves.png}
  \caption{CPDM training curves across both long-run launches combined. Top row: train loss, single-batch val loss, and per-epoch val loss. Bottom row: paper-aligned validation metrics (LPIPS, SSIM, PSNR) computed each epoch via the full 200-step Brownian Bridge reverse process on 16 val images. The narrow oscillation visible after epoch 1 motivates the manual stop at epoch 11.}
  \label{fig:cpdm-curves}
\end{figure}

\subsection{Baseline comparison}
\label{sec:baselines}

To verify that the model is doing more than predicting the data floor of a 
background-dominated distribution, CPDM was compared against three trivial 
baselines on a fixed sample of 200 random test slices using identical metric 
formulations:

\begin{table}[h]
\centering
\small
\begin{tabular}{lrrrr}
\toprule
Predictor & LPIPS$\downarrow$ & MAE$\downarrow$ & SSIM$\uparrow$ & PSNR$\uparrow$ \\
\midrule
Predict $-1$ everywhere      & 0.1767 & 0.0060 & 0.8945 & 37.45 \\
Predict the mean-PET image   & 0.1480 & 0.0065 & 0.9136 & 38.29 \\
Predict the CT (identity)    & 0.2427 & 0.0389 & 0.6869 & 24.32 \\
\textbf{CPDM (this work)}    & $\mathbf{\approx 0.15}$ & ${\approx}\,0.008$ & ${\approx}\,0.91$ & $\mathbf{\approx 40.6}$ \\
\bottomrule
\end{tabular}
\caption{CPDM vs.\ three trivial baselines on 200 randomly-sampled \texttt{autoPET} test slices.}
\label{tab:baselines}
\end{table}

Interpretation:
\begin{itemize}
  \item \textbf{PSNR}: CPDM beats every trivial baseline by 2--16\,dB --- the metric most sensitive to structural placement of bright regions.
  \item \textbf{LPIPS}: tied with the mean-PET baseline (both look ``plausibly PET'' to the perceptual network); clearly better than pure-black or identity.
  \item \textbf{SSIM}: tied --- dominated by the agreed-upon background in this dataset.
  \item \textbf{MAE}: CPDM \emph{loses} to mean-PET by ${\approx}24\%$, which is the expected signature of a model that produces actual variability (every misplaced hotspot costs MAE versus a degenerate mean predictor).
\end{itemize}

This baseline comparison confirms that the model is doing real structural work rather than collapsing to the data floor. The MAE-loss-with-PSNR-win combination is a textbook bias-variance pattern: the model invests its budget in variance-producing predictions, which PSNR rewards more than MAE.

\subsection{Qualitative results}

The most informative slice is \texttt{06\_cpdm\_samples.png}, a $6 \times 4$ 
grid showing for four held-out test slices: CT input, attention map, attenuation map, 
generated PET, ground-truth PET, and $L_1$ error.

\paragraph{Slice selection.} The four slices shown in 
Figure~\ref{fig:cpdm-samples} were \emph{not} drawn uniformly at random. 
An earlier version of this figure 
produced PET panels that looked almost completely black. The reason is a property 
of the dataset, not of the model: the brain-only crop concentrates roughly 80-90\% 
of pixels at SUV $\approx 0$ (background), and a uniformly random slice from this 
distribution is overwhelmingly likely to be background-dominated with no 
meaningful uptake to render. Visualising such a slice in the \texttt{hot} 
colormap correctly produces a near-black image because there is no signal to show. 
To give a faithful impression of what the model can actually produce, slices 
are instead ranked by their ground-truth PET maximum and four are selected: 
two from the top of the test distribution (strong-uptake case) and two from 
the middle (typical case). The bottom quartile (purely background slices) 
is excluded because a near-zero target makes the figure visually uninformative 
regardless of model output. Quantitatively, the full distribution including 
these low-uptake slices is reflected in the per-epoch metrics of 
Table~\ref{tab:cpdm-trajectory}, which are computed over a random 16-slice 
validation subset every epoch.

In addition, the \texttt{hot} colormap is stretched per column to the 99th percentile of 
the GT slice's display range (same \texttt{vmax} applied to GT and generated panels), 
so that brain-uptake intensities, which sit in the lower part of the SUV $[0,32]$ window, 
render as visibly bright rather than dark red. The same image with \texttt{vmax}=1.0 
looks substantially blacker, but that range is calibrated to lesion-level peak 
SUV values that almost never occur in brain slices.

Per-slice metrics on this batch:

\begin{table}[h]
\centering
\small
\begin{tabular}{rlrrrr}
\toprule
Slice & Uptake band & LPIPS$\downarrow$ & MAE$\downarrow$ & SSIM$\uparrow$ & PSNR$\uparrow$ \\
\midrule
369  & top of distribution    & 0.185 & 0.0138         & 0.896          & 27.23 \\
1948 & top of distribution    & 0.190 & 0.0158         & 0.879          & 28.85 \\
6137 & middle of distribution & 0.174 & 0.0104         & 0.889          & 35.48 \\
2037 & middle of distribution & \textbf{0.129} & \textbf{0.0077} & \textbf{0.917} & \textbf{42.46} \\
\bottomrule
\end{tabular}
\caption{Per-slice metrics on the qualitative batch rendered by \texttt{report.py}. Slices are ranked by ground-truth PET maximum and two are sampled from the top of the validation distribution (saturating slices with maximum normalised PET $= 1.0$) and two from the middle. The bottom quartile (background-only slices) is excluded as the figure would be uninformative there.}
\label{tab:qualitative}
\end{table}

\begin{figure}[h]
  \centering
  \includegraphics[width=\linewidth]{../report_out/06_cpdm_samples.png}
  \caption{CPDM end-to-end qualitative samples on four held-out test slices, selected by ground-truth PET maximum uptake (two from the top of the val distribution, two from the middle) so the panel is not dominated by background-heavy brain slices. Rows: CT input, attention map, CT-derived 511\,keV attenuation map, generated PET, ground-truth PET, and per-pixel $L_1$ error. PET is shown in the standard \texttt{hot} colormap (black\,$\to$\,red\,$\to$\,orange\,$\to$\,yellow\,$\to$\,white) with per-column \texttt{vmax} set to the 99th percentile of the GT slice's display values; the same \texttt{vmax} is applied to both the generated and the GT panel in a column so they are directly comparable. Without this stretch, typical brain-uptake values (SUV $\approx 5$--$15$) fall in the lower 15--45\% of the colormap and appear nearly black, which is a display artifact, not a model failure --- the model's quantitative fidelity is reported in Table~\ref{tab:cpdm-trajectory}.}
  \label{fig:cpdm-samples}
\end{figure}

\section{Discussion}

The reimplementation produces images that place uptake roughly where the 
attention map indicates, and it clears the trivial baselines on metrics 
that reward structural placement. On the absolute numbers, however, 
it does \emph{not} match the figures reported in the original CPDM paper. 
The reasons for that gap are concrete and identifiable:
\begin{enumerate}
  \item \textbf{Resolution.} At $64 \times 64$ the VQGAN bottleneck ($4 \times 16 \times 16$ latent, $4\times$ compression) caps the achievable detail. The paper operates at $256 \times 256$; doubling our resolution to $128 \times 128$ already recovers a meaningful fraction of that headroom.
  \item \textbf{Reduced VQGAN training objective.} The autoencoder was trained
  with only L1 reconstruction + codebook loss, omitting the perceptual 
  (LPIPS) and adversarial discriminator terms that VQGAN was originally 
  designed around. Those omitted terms are precisely what supply the 
  high-frequency texture supervision, so the decoder being useed is structurally 
  less expressive than the paper's.
  \item \textbf{Dataset scope.} The brain-only crop concentrates background pixels ($\sim 80$--$90\,\%$) and presents only a narrow slice of the variability the model would otherwise have to represent. Whole-body \texttt{autoPET}, as used in the paper, contains a far wider range of uptake patterns and anatomical context to learn from.
\end{enumerate}

None of these are methodological failures of CPDM itself. The pipeline is 
architecturally faithful to the paper. They are the cost of fitting the 
method into a single-CPU budget, and each is straightforward to lift in a 
follow-up with more compute and data.

The baseline comparison in Section~\ref{sec:baselines} is the appropriate 
yardstick for what this practical work can demonstrate: the model 
is better than data-floor predictors on the same data, even if it 
cannot be claimed to reach the original paper's absolute quality.

\section{Limitations and Future Work}

\begin{itemize}
  \item \textbf{CPU compute only}: no large-scale or high-resolution comparison possible.
  \item \textbf{Brain-only data}: limited evaluation domain.
  \item \textbf{No body crop in preprocessing.} Slices retain the full CT field of view including empty background around the patient. This inflates the near-zero-SUV pixel fraction and likely compresses the useful learning signal; a tight body-mask crop before resizing to $64 \times 64$ would concentrate more pixels on tissue and would plausibly improve the quality of the generated PET, but was not attempted here.
\end{itemize}

\section{Conclusion}

All four stages of CPDM (attention-map UNet, bulk export, VQGAN, Brownian Bridge 
diffusion) were implemented and trained end-to-end on \texttt{autoPET}. 
The pipeline runs from preprocessing to sampling, training and validation losses decrease and 
plateau rather than diverge, and the quantitative metrics are above 
the trivial baselines, so the model is learning something about the CT-to-PET 
mapping rather than fitting noise. It does not, however, reach the fidelity 
reported in the original CPDM paper. The generated images do not look like a 
faithful PET reconstruction either: in practice they resolve to a blurry red 
blob roughly in the centre of the slice, which is at least a step above an 
all-black output but is clearly not a usable PET image.

\printbibliography

\end{document}
